\documentclass[11pt]{article}
\usepackage[margin=1in]{geometry}
\usepackage{booktabs}
\usepackage{array}
\usepackage{amsmath}
\usepackage{url}

% This is a static presentation copy generated from docs/permit_types_brief.md,
% which is the canonical, actively-edited source. It is NOT machine-generated by
% a script and has no \input fragments -- if the .md changes, regenerate this
% file by hand (or ask Claude to) rather than editing the two out of sync.

\title{\vspace{-2.5em}NPDES Permitting}
\author{}
\date{\today}

\begin{document}
\sloppy
\maketitle
\vspace{-2.5em}

\section*{Introduction}
This brief covers the scope of National Pollutant Discharge Elimination System
(NPDES) permitting under the Clean Water Act (CWA), how a permit is classified
within that scope (individual vs.\ general, and major vs.\ minor), and how those
classifications show up in the ICIS-NPDES (ECHO) data this project uses. Also
covers the Master General Permits download, enforcement patterns by permit type,
the 2015 Electronic Reporting Rule, and why many projects restrict to major
individual permits.

\section{Scope}

NPDES (CWA \S\,402) requires a permit for the discharge of a pollutant from a
\textbf{point source} into waters of the United States. A point source is defined
as any discernible, confined and discrete conveyance, including pipes, ditches, and
vessels. Several major discharge categories are not regulated under NPDES:

\begin{itemize}
  \item \textbf{Nonpoint source pollution} --- runoff not channeled through a
        discrete conveyance (general agricultural field runoff, most urban/lawn
        runoff not routed through a regulated storm sewer, atmospheric deposition).
        It is addressed through CWA \S\,319 nonpoint source management grants,
        which have no permitting mechanism.
  \item \textbf{Agricultural stormwater and irrigation return flows} --- excluded
        from the statutory definition of ``point source'' itself (CWA
        \S\,502(14)), even when somewhat channeled (e.g.\ field drainage). Animal
        feeding operations are not permitted under NPDES, but facilities
        designated as concentrated animal feeding operations are permitted.
  \item \textbf{Silvicultural (forestry) nonpoint runoff} --- ordinary forestry
        runoff is excluded, though certain forest-road discharges through a
        discrete conveyance can still count as a point source.
  \item \textbf{Dredge-and-fill activities} --- regulated under a different CWA
        provision (\S\,404, Army Corps of Engineers), not \S\,402/NPDES.
  \item \textbf{Zero-discharge facilities} --- a facility with no discharge to a
        jurisdictional water (pure land application, fully contained operations)
        has nothing to permit, since there's no point-source discharge in the
        first place.
\end{itemize}

\paragraph{Contested boundaries.}
\begin{itemize}
  \item \textbf{Groundwater} --- historically treated as largely outside NPDES,
        but \textit{County of Maui v.\ Hawaii Wildlife Fund} (2020) held a
        discharge reaching surface water via groundwater can require a permit if
        it is the ``functional equivalent of a direct discharge'' --- a
        fact-specific test, not a bright line.
  \item \textbf{``Waters of the United States'' (WOTUS)} --- which water bodies
        count as jurisdictional at all has been repeatedly redefined by rule and
        litigation across administrations; a discharge to a water outside the
        current WOTUS definition is outside NPDES regardless of point-source
        status.
\end{itemize}

\section{Individual vs.\ General Permits}

Within that scope, the CWA provides two vehicles for issuing an NPDES permit:

\bigskip
An \textbf{individual permit} is written for one specific facility. Its effluent
limits combine technology-based limits (national, by industry) with
\textbf{water-quality-based limits derived from the specific receiving water}, so
two identical dischargers can face different limits. It requires a full,
facility-specific application (40 CFR \S\,122.21) and carries facility-specific
monitoring and reporting (DMR) obligations. Individual permits are for larger,
more complex dischargers, such as municipal treatment plants and major industrial
sources.

\bigskip
A \textbf{general permit} (40 CFR \S\,122.28) is a single template permit the
permitting authority issues once to cover \textbf{many similar, lower-risk
dischargers} at once (construction and industrial stormwater, CAFOs, small or
short-lived operations). A facility obtains coverage by filing a \textbf{Notice of
Intent (NOI)} against the master permit. Under general permits, conditions are
standardized and monitoring is lighter or absent; one master general permit can
cover hundreds of thousands of individual coverages (Section~\ref{sec:dataset}).

\bigskip
\paragraph{Permit term:} Both permit types are capped at the same statutory
maximum --- 40 CFR \S\,122.46 limits any NPDES permit, individual or general, to a
fixed term not exceeding five years. Reissuing one general permit covers every
facility under it at once, so general permits are almost always issued for the
full five years, while an individual permit's specific term (and whether it's
reissued on time or instead runs on under administrative continuance) depends on
that facility's own permitting history.

\paragraph{Representative general-permit programs:} EPA's Multi-Sector General
Permit (MSGP) covers stormwater discharges from industrial facilities across
roughly thirty sectors in areas where EPA is the permitting authority; the
Construction General Permit (CGP) covers stormwater from construction sites
disturbing $\geq$1 acre with a pathway to a water of the US. CAFOs are typically
covered under state- or EPA-issued general permits as well.

The practical trade-off is specificity vs.\ administrative cost: individual
permits are tailored but resource-intensive; general permits are efficient but
coarse.

\section{Facility size: major vs.\ minor}
\label{sec:majorminor}

Independent of \textit{how} a permit is issued (Section~2), every permit also
carries a \textbf{major/minor} designation.

EPA's major/minor designation comes from a point-based rating system (originally
the Municipal/Industrial Strategic Initiative rating) that scores factors like
discharge flow volume, toxic-pollutant potential, and the public-health and
water-quality impact of the receiving water. Facilities crossing the point
threshold are designated \textbf{major}. For municipal treatment plants,
facilities with a design flow of roughly $\geq$1 million gallons per day are
typically treated as major. In practice, this means closer EPA/state oversight,
more frequent inspections, and --- per Section~6 --- the primary population where
formal enforcement and SNC tracking concentrate. Major status is re-evaluated
over a permit's life and can change between reissuances (hence ``ever major,''
used throughout this brief, rather than a permanent label).

\paragraph{DMR reporting frequency} is where major/minor status shows up most
directly in the data. Major permits typically require \textbf{monthly} numeric
effluent monitoring and DMR submission across every permitted parameter, since
majors carry the facility-specific technology-based and water-quality-based
limits described in Section~2. Minor permits usually specify a much lighter
cadence --- quarterly, semi-annual, or even annual --- and many minor facilities
covered under a general permit have no numeric DMR obligation at all, submitting
only a periodic certification of best-management-practice compliance rather than
measured effluent values. This reporting-frequency gap is the direct mechanical
cause of the $\sim$84\% vs $\sim$7\% DMR coverage split already noted in
Section~\ref{sec:dataset}, and it's also why the 2015 Electronic Reporting Rule's
phased rollout (Section~7) treated individual/major data as Phase 1 and
general/minor data as Phase 2: there was simply far less minor-facility DMR data
to migrate to electronic reporting in the first place.

\paragraph{Major and minor cut across both permit vehicles, asymmetrically.}

\begin{table}[h]
\centering
\caption{Major/minor status by permit vehicle (distinct permits, ``ever''
flagged in any version).}\label{tab:majorminor}
\begin{tabular}{l r r r r}
\toprule
Vehicle & Total permits & Ever major & Ever minor & \% ever major \\
\midrule
\texttt{NPD} (individual) & 99,940 & 7,800 & 92,140 & 7.8\% \\
\texttt{GPC} (general) & 1,029,666 & 95 & 1,029,571 & 0.009\% \\
Other codes (\texttt{UFT}/\texttt{IIU}/\texttt{APR}/\texttt{SIN}/\texttt{NGP}/\texttt{SNN}) & 64,424 & 167 & 64,257 & 0.26\% \\
\bottomrule
\end{tabular}
\end{table}

\textbf{Major status is common among individual permits (7.8\%) and vanishingly
rare among general permits (0.009\%)} --- it's not just that most majors happen
to be individual, it's that being individually permitted in the first place is
strongly associated with being major, and the reverse is nearly never true for
general-permit coverages.

\section{Distinction in the dataset}
\label{sec:dataset}

\paragraph{Permit-type codes:} \texttt{PERMIT\_TYPE\_CODE} marks the vehicle:
\texttt{NPD} = individual, \texttt{GPC} = a facility's coverage under a general
permit, \texttt{NGP} = the general permit itself (the master record, not a
covered facility). The remaining codes (\texttt{UFT}, \texttt{IIU},
\texttt{APR}, \texttt{SIN}, \texttt{SNN}) are minor, low-volume variants
(state-issued or non-NPDES-equivalent records).

\begin{table}[h]
\centering
\caption{\texttt{PERMIT\_TYPE\_CODE} frequency. Permits = distinct
\texttt{EXTERNAL\_PERMIT\_NMBR} values; Records = raw rows in
\texttt{ICIS\_PERMITS.csv} (one row per permit version).}\label{tab:codes}
\begin{tabular}{l r r}
\toprule
Code & Permits (distinct) & Records (all versions) \\
\midrule
\texttt{GPC} & 1,029,666 & 1,368,680 \\
\texttt{NPD} & 99,940 & 242,678 \\
\texttt{UFT} & 46,360 & 46,360 \\
\texttt{IIU} & 7,315 & 11,105 \\
\texttt{APR} & 5,901 & 16,382 \\
\texttt{SIN} & 3,573 & 6,562 \\
\texttt{NGP} & 1,232 & 2,823 \\
\texttt{SNN} & 41 & 56 \\
\bottomrule
\end{tabular}
\end{table}

The gap matters: general-permit coverages (\texttt{GPC}) outnumber individual
permits roughly ten to one on either count, but reissuance behavior differs ---
\textbf{51.9\%} of individual permits carry more than one version (a permit keeps
its \texttt{NPDES\_ID} across $\sim$5-year reissuances), versus \textbf{21.8\%}
of general-permit coverages, which are more often \textbf{renumbered} when the
master permit is reissued (e.g.\ an \texttt{NDR100000}-cycle coverage becomes an
\texttt{NDR110000}-cycle coverage under a new ID) rather than versioned under the
same one.

\begin{itemize}
  \item \textbf{Facility identity:} Individual permits map cleanly to a single
        physical site with a reliable FRS registry ID (\texttt{FACILITY\_UIN})
        and point geocode. General permits frequently \textbf{share one generic
        \texttt{FACILITY\_UIN}} across unrelated sites. This makes
        facility-level aggregation reliable for individual permits and
        hazardous for general ones.
  \item \textbf{Multiple permits:} 10,362 \texttt{FACILITY\_UIN}s hold both an
        individual (\texttt{NPD}) and a general (\texttt{GPC}) permit ---
        typically an individually permitted process discharge and a general
        stormwater coverage at the same site. The two permits govern different
        discharges, not the same one, so they shouldn't be collapsed into a
        single ``one facility, one permit'' record.
  \item \textbf{Data richness:} Monitoring and enforcement are concentrated on
        the individual/major side. In FY2025 DMR data, majors appear at
        $\sim$84\% reporting coverage vs.\ $\sim$7\% for minors --- consistent
        with general permits' lighter-to-absent numeric monitoring
        requirements.
\end{itemize}

\section{The Master General Permits download}

The 2,823 \texttt{NGP} master records don't come from the main 15-table
ICIS-NPDES bundle (\path{npdes_downloads.zip}) --- they have their own
\textbf{separate} download: \path{npdes_master_general_permits.zip}, fetched from
\url{https://echo.epa.gov/files/echodownloads/npdes_master_general_permits.zip}.

\begin{itemize}
  \item \textbf{One row per master permit} (2,823 total, all
        \texttt{PERMIT\_TYPE\_CODE = "NGP"}) --- this file is \textit{only} the
        master templates, not the millions of individual coverages.
  \item \textbf{Extra descriptive fields not needed for a coverage record} ---
        \texttt{PERMIT\_NAME} (e.g.\ ``Hawaii Master General Permit Noncontact
        Cooling Waters''), \texttt{ISSUING\_AGENCY},
        \texttt{ORIGINAL\_ISSUE\_DATE}/\texttt{ISSUE\_DATE}/%
        \texttt{EFFECTIVE\_DATE}/\texttt{EXPIRATION\_DATE}, and the same
        \texttt{EXTERNAL\_PERMIT\_NMBR} that appears as the target of
        \texttt{MASTER\_EXTERNAL\_PERMIT\_NMBR} on \texttt{GPC} records in the
        main \texttt{ICIS\_PERMITS.csv}.
  \item \textbf{\texttt{MASTER\_EXTERNAL\_PERMIT\_NMBR} is (almost) always blank
        in this file itself} --- 51 of 2,823 rows are an exception and have it
        populated, worth a closer look before treating ``blank = top-level
        master'' as an absolute rule.
\end{itemize}

\section{Enforcement patterns by permit type}

Formal vs.\ informal enforcement action records, by permit type and major/minor
status (all years, all actions in the bulk enforcement files; each row is one
action record, not one facility):

\begin{table}[h]
\centering
\footnotesize
\caption{Formal vs.\ informal enforcement action records.}\label{tab:enforcement}
\begin{tabular}{l l r r r r}
\toprule
Permit type & Status & Formal & Informal & Total & \% Informal \\
\midrule
Individual & Minor & 42,307 & 351,145 & 393,452 & 89.2\% \\
General & Minor & 23,313 & 237,934 & 261,247 & 91.1\% \\
Individual & Major & 30,942 & 179,192 & 210,134 & 85.3\% \\
Other & Unknown & 12,748 & 27,885 & 40,633 & 68.6\% \\
Other & Minor & 1,406 & 24,147 & 25,553 & 94.5\% \\
General & Unknown & 513 & 427 & 940 & 45.4\% \\
Other & Major & 279 & 526 & 805 & 65.3\% \\
General & Major & 59 & 596 & 655 & 91.0\% \\
Individual & Unknown & 249 & 125 & 374 & 33.4\% \\
\bottomrule
\end{tabular}
\end{table}

Pooling across major/minor status, individual permits accrue \textbf{603,960}
action records (12.2\% formal / 87.8\% informal) and general permits
\textbf{262,842} (9.1\% formal / 90.9\% informal). The two permit types look
broadly similar in their formal/informal \textit{split}, but arrive there by
different detection pathways:

\begin{itemize}
  \item \textbf{Individual permits} (mostly majors) generate violations chiefly
        through routine DMR review (numeric effluent-limit exceedances) and
        scheduled inspections. Individual/major permittees are also the
        population to which EPA's \textbf{significant noncompliance (SNC)}
        criteria are principally applied: since 1995, EPA policy has defined SNC
        to include, among other triggers, any monthly average effluent limit
        exceeded by $\geq$40\% (or by $\geq$20\% for two or more consecutive
        months), any limit violated for four or more months in a
        two-consecutive-quarter window, unauthorized bypass/pass-through, and
        reporting violations 30+ days late. These criteria presuppose numeric
        monthly DMR limits --- exactly what individual permits carry and general
        permits typically do not --- which is why SNC tracking is substantively
        an individual/major-permit concept rather than a general-permit one.
  \item \textbf{General permits} have less routine numeric monitoring, so
        violations more often surface through inspections, failure to file the
        NOI or other required paperwork, or complaints, rather than a DMR
        exceedance. Consistent with this, the response is more often informal
        (90.9\% vs.\ 87.8\% for individual). The Director's authority to revoke
        general-permit coverage and require an individual permit (Section~2; 40
        CFR \S\,122.28(b)(3)) is the formal escalation path when a covered
        facility turns out to warrant closer, facility-specific regulation.
\end{itemize}

\section{The 2015 NPDES Electronic Reporting Rule}

EPA's NPDES Electronic Reporting Rule took effect \textbf{December 21, 2015},
replacing paper-based NPDES reporting with mandatory electronic submission, with
the stated goal of more complete, timely, and nationally consistent compliance
data. Implementation was phased: \textbf{Phase 1} (largely individual-permit DMR
data) had a compliance deadline of December 21, 2016; \textbf{Phase 2}
(general-permit and other remaining data streams) was originally due December
21, 2020 but was later extended by rule, giving affected programs up to five
additional years. This phased rollout is the mechanism behind the ``2016 eRule''
reporting break already flagged in this project's data documentation
(\path{docs/data_quirks.md}): majors' DMR coverage was already high before the
rule, while minors' and general-permit coverage only began rising as Phase 1,
and later Phase 2, requirements took hold --- so pre- and post-2016 (and
pre/post-2020+) coverage are not directly comparable, particularly for
general-permit-covered dischargers.

\section{Implications for analysis}

For facility-level research, the individual permit is the meaningful analytical
unit: one physical site, a stable identifier, dense measurement, and a
trustworthy location. The general permit is best understood as an
\textbf{administrative umbrella, not a facility} --- a single
\texttt{FACILITY\_UIN} or master permit number can span hundreds of unrelated
sites, and monitoring/enforcement data is comparatively sparse. Restricting to
major individual permits, as most academic projects do, produces a well-defined
and well-measured population. However, it excludes the large majority of
permitted dischargers, which are general, lightly monitored, and (per
Section~6) usually resolved informally rather than through the
DMR/inspection/SNC pipeline that individual permits generate.

\section*{Reproducibility}
This document is a static presentation copy of \path{docs/permit_types_brief.md}
(the actively-edited, canonical source --- edit that file, not this one).
Figures were computed from \path{data/raw/npdes_downloads/ICIS_PERMITS.csv},
\path{data/raw/npdes_downloads/ICIS_FACILITIES.csv}, the separately-downloaded
\texttt{data/raw/Master General Permits/ICIS\_MASTER\_GENERAL\_PERMITS.csv}, and
the pre-built \path{output/enforcement_by_permit_type.csv} (from
\path{code/diagnostics/enforcement_breakdowns/enforcement_by_permit_type.R});
rounded. Regulatory citations (40 CFR \S\S\,122.21, 122.28, 122.46; the 1995 SNC
criteria memorandum; the 2015 Electronic Reporting Rule and its Phase 2
extension) are drawn from the Federal Register and eCFR and are not derived from
project data.

\end{document}
